% !PRINTED 501
\weberpage{501}
%<<CONT
Wert hat, so ist sie eine ganze Funktion von $j(\omega)$. Lassen
wir in (5) $x$ in $v$ übergehen, so ergibt sich
\weq{6}{M = \frac{\psi(v,u)}{F'(v)},}
worin $\psi(v,u)$ eine ganze rationale Funktion von $u$ und $v$
ist, mit rationalen Zahlenkoeffizienten.

Lassen wir nun in (6) $\omega$ in die Wurzel einer quadratischen
Gleichung der negativen Diskriminante
\[
D = B^2 - 4AC
\]
übergehen:
\weq{7}{A\omega^2 + B\omega + C = 0,}
so können einer oder mehrere der Werte $v$ gleich $u$ werden. Ist
dies nur für einen der $\nu$ Werte $v$ der Fall, so wird $F'(v)$ für
$v = u$ nicht verschwinden, und wir können nach (6) $M$ rational
durch diesen singulären Wert von $u$ ausdrücken.

Die Gleichung $u = v$ führt aber die Bedingung mit sich:
\[
\frac{c+\partial\omega}{a} = \frac{\gamma+\delta\omega}{\alpha+\beta\omega}.
\]
Daraus folgt:
\weq{8}{\begin{aligned}
\partial\beta &= Ax, & \partial\alpha &= \frac{Bx-y}{2},\\
c\beta-a\delta &= \frac{Bx+y}{2}, & c\alpha-a\gamma &= Cx,
\end{aligned}}
\weq{9}{4n = y^2 - Dx^2,}
worin $x$ positiv angenommen werden kann.

Wir nehmen $A$ ungerade und relativ prim zu $D$ an. Ist
$D \equiv 1\ (\mathrm{mod}\ 4)$, so setzen wir $n = -D$ und erhalten aus (9):
\[
y^2 + n(x^2-4) = 0.
\]
Der Wert $x=1$ ist hiernach nur dann zulässig, wenn $-D$
das Dreifache einer Quadratzahl ist. Sehen wir zunächst von
diesem Fall ab, so bleibt nur
\[
x = 2, \quad y = 0,
\]
und aus (8) ergibt sich:
\[
\begin{aligned}
\partial &= 1, & \beta &= 2A, & \alpha &= B,\\
a &= n, & \gamma &\equiv 2C, & \delta &\equiv B\ (\mathrm{mod}\ 4),\\
\alpha + \beta\omega &= \sqrt{D}, &&&&
\end{aligned}
\]
und durch die beiden letzten Gleichungen (8) ist $c$ nach dem
Modul $n$ bestimmt.
